\chapter[Lezione 11 --- 21 ottobre 2026]{Lezione 11: intorni e proprietà dei limiti}
\label{chap:lezione-11}

\section{Intersezioni di intorni}
\label{sec:lez11-intersezione-intorni}

Sia \(x_0\in\overline{\R}\). L'intersezione di due intorni di
\(x_0\) è ancora un intorno di \(x_0\).

Se \(x_0\in\R\), per la \cref{def:lez07-intorno} si ha, per
ogni \(\delta_1,\delta_2>0\),
\[
  D(x_0,\delta_1)\cap D(x_0,\delta_2)
  =D\bigl(x_0,\min\{\delta_1,\delta_2\}\bigr).
\]
Se \(x_0=+\infty\), per la \cref{def:lez07-intorni-infinito},
per ogni \(k_1,k_2\in\R\),
\[
  (k_1,+\infty)\cap(k_2,+\infty)
  =(\max\{k_1,k_2\},+\infty).
\]
Analogamente, se \(x_0=-\infty\),
\[
  (-\infty,k_1)\cap(-\infty,k_2)
  =(-\infty,\min\{k_1,k_2\}).
\]
In tutti i casi, l'intersezione è un intorno del medesimo
punto. Applicando ripetutamente la stessa proprietà, si
ottiene il risultato per l'intersezione di un numero finito di intorni di \(x_0\).

\begin{figure}[htbp]
  \centering
  \begin{tikzpicture}[unipd/diagram]
  \draw[->,unipd-muted] (-.6,-1.05) -- (8,-1.05) node[right] {\(x\)};
  \draw[very thick,unipd-primary,->] (1.45,.85) -- (7.55,.85);
  \draw[unipd-primary,fill=white,thick] (1.45,.85) circle (2.4pt);
  \draw[very thick,unipd-accent,->] (3.45,.10) -- (7.55,.10);
  \draw[unipd-accent,fill=white,thick] (3.45,.10) circle (2.4pt);
  \draw[very thick,unipd-primary,->] (3.45,-.60) -- (7.55,-.60);
  \draw[unipd-primary,fill=white,thick] (3.45,-.60) circle (2.4pt);
  \draw[dashed,unipd-muted] (1.45,-1.13) -- (1.45,1.08);
  \draw[dashed,unipd-muted] (3.45,-1.13) -- (3.45,1.08);
  \node[below=6pt] at (1.45,-1.05) {\(k_1\)};
  \node[below=6pt] at (3.45,-1.05) {\(k_2\)};
  \node[above=4pt,text=unipd-primary] at (5.7,.85) {\((k_1,+\infty)\)};
  \node[above=4pt,text=unipd-accent] at (5.7,.10) {\((k_2,+\infty)\)};
  \node[above=4pt,text=unipd-primary] at (5.7,-.60) {\((\max\{k_1,k_2\},+\infty)\)};
\end{tikzpicture}

  \caption{Intersezione di due intorni di \(+\infty\):
  la semiretta comune inizia dalla maggiore delle due soglie,
  qui con \(k_1<k_2\).}
  \label{fig:lez11-intersezione-intorni-infinito}
\end{figure}

L'affermazione non si estende, in generale, alle intersezioni
infinite, come mostra l'esempio seguente.

\begin{example}[Intersezione di infiniti intorni]
\label{ex:lez11-intersezione-infinita}
Per ogni \(n\in\N\) con \(n\ge1\), si consideri
\[
  A_n:=D\!\left(0,\frac1n\right)
       =\left(-\frac1n,\frac1n\right).
\]
Allora
\[
  B:=\bigcap_{n=1}^{\infty}A_n=\{0\}.
\]
L'insieme \(B\) non è quindi un intorno di \(0\), benché
ciascun \(A_n\) lo sia.
\end{example}

\begin{proof}
Poiché \(0\in A_n\) per ogni \(n\ge1\), si ha
\(0\in B\).

Si consideri l'insieme
\[
  E:=\left\{\frac1n:n\in\N,\ n\ge1\right\}.
\]
Il numero \(0\) è un minorante di \(E\). Inoltre, per
ogni \(\varepsilon>0\), per il \cref{thm:proprieta-archimedea}
esiste \(n\in\N\) tale che \(n>1/\varepsilon\);
pertanto \(1/n<\varepsilon\). Per il
\cref{thm:caratterizzazione-inf}, segue
\[
  \inf E=0.
\]
Dall'\cref{ex:estremi-insieme-opposto} si ricava anche
\[
  \sup(-E)=-\inf E=0.
\]

Sia ora \(x\in B\). Per definizione di intersezione,
\[
  -\frac1n<x<\frac1n
  \qquad\forall n\in\N,\ n\ge1.
\]
Il numero \(x\) è dunque un minorante di \(E\) e un
maggiorante di \(-E\). Ne segue che
\[
  x\le\inf E=0,
  \qquad
  x\ge\sup(-E)=0.
\]
Pertanto \(x=0\), e quindi \(B\subseteq\{0\}\).
Insieme all'inclusione opposta, si conclude
\[
  B=\{0\}.
\]
\end{proof}

\section{Unicità del limite}
\label{sec:lez11-unicita-limite}

La \cref{def:lez08-limite} non permette a una funzione
di ammettere due limiti distinti nello stesso punto di
accumulazione.

\begin{theorem}[Unicità del limite]
\label{thm:lez11-unicita-limite}
Sia \(f\colon A\longrightarrow\R\), con \(A\subseteq\R\)
non vuoto, e sia \(x_0\in\overline{\R}\) un punto di
accumulazione per \(A\). Se
\[
  \lim_{x\to x_0}f(x)=\ell_1,
  \qquad
  \lim_{x\to x_0}f(x)=\ell_2,
\]
con \(\ell_1,\ell_2\in\overline{\R}\), allora
\[
  \ell_1=\ell_2.
\]
\end{theorem}

\begin{proof}
Si procede per assurdo, supponendo
\(\ell_1\ne\ell_2\).

Esistono due intorni \(V_{\ell_1}\) e \(V_{\ell_2}\) dei
rispettivi punti tali che
\[
  V_{\ell_1}\cap V_{\ell_2}=\varnothing.
\]
Infatti, se \(\ell_1,\ell_2\in\R\), si può supporre,
senza perdita di generalità, \(\ell_1<\ell_2\) e porre
\[
  \delta:=\frac{\ell_2-\ell_1}{2}>0,
  \qquad
  V_{\ell_1}:=D(\ell_1,\delta),
  \qquad
  V_{\ell_2}:=D(\ell_2,\delta).
\]
I due intorni sono disgiunti, come illustrato nella
\cref{fig:lez11-intorni-disgiunti}. Se almeno uno dei
due limiti è infinito, si possono scegliere intorni
disgiunti utilizzando la \cref{def:lez07-intorni-infinito}.

Per la \cref{def:lez08-limite}, esistono intorni
\(U_{x_0}\) e \(U'_{x_0}\) di \(x_0\) tali che
\[
  f(x)\in V_{\ell_1}
  \quad\forall x\in(U_{x_0}\cap A)\setminus\{x_0\},
\]
\[
  f(x)\in V_{\ell_2}
  \quad\forall x\in(U'_{x_0}\cap A)\setminus\{x_0\}.
\]
Per la \cref{sec:lez11-intersezione-intorni},
\(W_{x_0}:=U_{x_0}\cap U'_{x_0}\) è ancora un intorno
di \(x_0\). Poiché \(x_0\) è di accumulazione per
\(A\), esiste
\[
  x\in(W_{x_0}\cap A)\setminus\{x_0\}.
\]
Per tale \(x\), entrambe le condizioni sul limite
devono essere verificate, quindi
\[
  f(x)\in V_{\ell_1}\cap V_{\ell_2}=\varnothing,
\]
una contraddizione. Pertanto \(\ell_1=\ell_2\).
\end{proof}

\begin{figure}[htbp]
  \centering
  \begin{tikzpicture}[unipd/diagram,scale=1.1]
  \draw[->,unipd-muted] (-.45,0) -- (8.1,0) node[right] {\(x\)};
  \draw[very thick,unipd-primary] (.7,0) -- (3.5,0);
  \draw[very thick,unipd-accent] (3.5,0) -- (6.3,0);
  \foreach \x in {.7,3.5} {
    \draw[unipd-primary,fill=white,thick] (\x,0) circle (2.4pt);
  }
  \foreach \x in {3.5,6.3} {
    \draw[unipd-accent,fill=white,thick] (\x,0) circle (2.0pt);
  }
  \fill[unipd-primary] (2.1,0) circle (2.1pt);
  \fill[unipd-accent] (4.9,0) circle (2.1pt);
  \node[below=7pt] at (2.1,0) {\(\ell_1\)};
  \node[below=7pt] at (4.9,0) {\(\ell_2\)};
  \node[above=8pt,text=unipd-primary] at (2.1,0) {\(V_{\ell_1}\)};
  \node[above=8pt,text=unipd-accent] at (4.9,0) {\(V_{\ell_2}\)};
  \draw[<->,unipd-muted] (2.1,.7) -- (3.5,.7) node[midway,above] {\(\delta\)};
  \draw[<->,unipd-muted] (3.5,.7) -- (4.9,.7) node[midway,above] {\(\delta\)};
  \draw[dashed,unipd-muted] (3.5,.12) -- (3.5,.85);
\end{tikzpicture}

  \caption{Due intorni circolari disgiunti di
  \(\ell_1<\ell_2\), con raggio
  \(\delta=(\ell_2-\ell_1)/2\).}
  \label{fig:lez11-intorni-disgiunti}
\end{figure}

\begin{exercise}[Intorni disgiunti nella retta reale estesa]
\label{ex:lez11-intorni-disgiunti}
Verificare che si possono scegliere due intorni disgiunti
anche quando \(\ell_1\ne\ell_2\) e almeno uno dei due
appartiene a \(\{-\infty,+\infty\}\), distinguendo
il caso in cui l'altro limite è reale da quello
in cui i due limiti sono infiniti.
\end{exercise}

\begin{proof}[Soluzione dell'esercizio]
Si distinguono tre casi.

\emph{Primo caso: \(\ell_1=+\infty\), \(\ell_2\in\R\).}
Si scelgano gli intorni
\[
  V_{\ell_1}:=(\ell_2+1,+\infty),
  \qquad
  V_{\ell_2}:=D(\ell_2,1)
  =(\ell_2-1,\ell_2+1).
\]
Poiché gli intervalli non hanno punti in comune,
si ha \(V_{\ell_1}\cap V_{\ell_2}=\varnothing\).

\emph{Secondo caso: \(\ell_1\in\R\), \(\ell_2=-\infty\).}
Si scelgano gli intorni
\[
  V_{\ell_1}:=D(\ell_1,1)
  =(\ell_1-1,\ell_1+1),
  \qquad
  V_{\ell_2}:=(-\infty,\ell_1-1).
\]
Anche in questo caso
\(V_{\ell_1}\cap V_{\ell_2}=\varnothing\).

\emph{Terzo caso: \(\ell_1=-\infty\), \(\ell_2=+\infty\).}
Si scelgano gli intorni
\[
  V_{\ell_1}:=(-\infty,0),
  \qquad
  V_{\ell_2}:=(0,+\infty).
\]
Si ha nuovamente
\(V_{\ell_1}\cap V_{\ell_2}=\varnothing\).

Scambiando gli indici si ottengono i casi rimanenti.
Pertanto, per ogni coppia di punti distinti della
retta reale estesa, esistono due intorni disgiunti.
\end{proof}

\section{Permanenza del segno}
\label{sec:lez11-permanenza-segno}

\begin{theorem}[Permanenza del segno]
\label{thm:lez11-permanenza-segno}
Sia \(f\colon A\longrightarrow\R\), con \(A\subseteq\R\)
non vuoto, e sia \(x_0\in\overline{\R}\) un punto di
accumulazione per \(A\).
\begin{enumerate}
  \item Se
  \[
    \lim_{x\to x_0}f(x)=\ell
    \in(0,+\infty)\cup\{+\infty\},
  \]
  allora esistono \(k>0\) e un intorno \(U_{x_0}\)
  di \(x_0\) tali che
  \[
    f(x)>k
    \qquad\forall x\in(U_{x_0}\cap A)\setminus\{x_0\}.
  \]
  \item Se
  \[
    \lim_{x\to x_0}f(x)=\ell
    \in(-\infty,0)\cup\{-\infty\},
  \]
  allora esistono \(k>0\) e un intorno \(U_{x_0}\)
  di \(x_0\) tali che
  \[
    f(x)<-k
    \qquad\forall x\in(U_{x_0}\cap A)\setminus\{x_0\}.
  \]
\end{enumerate}
\end{theorem}

\begin{proof}[Dimostrazione del primo punto]
Si distinguono i casi \(\ell=+\infty\) e
\(\ell\in(0,+\infty)\).

\emph{Caso \(\ell=+\infty\).}
Per la \cref{def:lez08-limite}, per ogni \(K\in\R\)
esiste un intorno \(U_{x_0}\) tale che
\[
  f(x)>K
  \qquad\forall x\in(U_{x_0}\cap A)\setminus\{x_0\}.
\]
Scegliendo \(K=1\), si ottiene la tesi con \(k=1\).

\emph{Caso \(\ell\in(0,+\infty)\).}
Per la \cref{def:lez08-limite}, per ogni
\(\varepsilon>0\) esiste un intorno \(U_{x_0}\)
tale che
\[
  |f(x)-\ell|<\varepsilon
  \qquad\forall x\in(U_{x_0}\cap A)\setminus\{x_0\}.
\]
Scegliendo \(\varepsilon=\ell/2>0\), segue
\[
  -\frac\ell2<f(x)-\ell<\frac\ell2,
\]
e quindi
\[
  f(x)>\frac\ell2>0
  \qquad\forall x\in(U_{x_0}\cap A)\setminus\{x_0\}.
\]
La tesi è verificata ponendo \(k=\ell/2\).
\end{proof}

Il secondo punto si dimostra analogamente, considerando
\(\ell=-\infty\) e \(\ell\in(-\infty,0)\); la verifica
non viene sviluppata in questa sede.

\section{Limitatezza locale}
\label{sec:lez11-limitatezza-locale}

\begin{theorem}[Limitatezza locale]
\label{thm:lez11-limitatezza-locale}
Sia \(f\colon A\longrightarrow\R\), con \(A\subseteq\R\)
non vuoto, e sia \(x_0\in\overline{\R}\) un punto di
accumulazione per \(A\). Se
\[
  \lim_{x\to x_0}f(x)=\ell\in\R,
\]
allora esiste un intorno \(U_{x_0}\) di \(x_0\)
tale che l'insieme
\[
  \{f(x):x\in U_{x_0}\cap A\}
\]
è limitato.
\end{theorem}

La dimostrazione del \cref{thm:lez11-limitatezza-locale}
non viene sviluppata in questa sede.
